\documentclass{article}
\usepackage{amsmath}
\usepackage{amssymb}
\usepackage{geometry}
\usepackage{graphicx}
\usepackage{xepersian}
\settextfont{XB Niloofar}

\geometry{a4paper, margin=1in}

\title{\lr{Kernel Regression: From Linear to Nonlinear Modeling} \\ (رگرسیون هسته: از مدل‌سازی خطی به غیرخطی)}
\author{}
\date{}
\begin{figure}[h!]
  \centering
  \includegraphics[width=0.75\textwidth]{img_page_7_1.png}
  \caption{نتایج برازش رگرسیون غیرخطی با استفاده از هسته گوسی}
\end{figure}

\begin{figure}[h!]
  \centering
  \includegraphics[width=0.75\textwidth]{img_page_8_1.png}
  \caption{تاثیر تغییرات منظم‌سازی و سطح نویز بر روی هموارسازی رگرسیون هسته}
\end{figure}






\begin{document}

\maketitle

\section*{فهرست مطالب}
\begin{enumerate}
    \item پایه‌ی رگرسیون خطی (\lr{Linear Regression})
    \item رگرسیون هسته (\lr{Kernel Regression}): ایده اصلی
    \item آموزش: حل برای $\alpha$
    \item پیش‌بینی روی داده‌های جدید
    \item ارتباط با شبه‌معکوس (\lr{Pseudoinverse})
    \item تاثیرات منظم‌سازی (\lr{regularization}) و سطوح نویز بر رگرسیون هسته
    \item دیدگاه نخستین (\lr{Primal}) در مقابل دوگان (\lr{Dual})
    \item نکات کلیدی
\end{enumerate}

\newpage

\section*{۱. پایه‌ی رگرسیون خطی}
هدف رگرسیون پیش‌بینی یک خروجی پیوسته $y$ از ویژگی‌های ورودی $x=(x_1, x_2, \dots, x_d)^T$ با استفاده از یک تابع یادگرفته‌شده $f$ است:
\[
y = f(x) + \varepsilon
\]
که در آن $\varepsilon$ نشان‌دهنده نویز تصادفی است.

در رگرسیون خطی (\lr{linear regression})، ما فرض می‌کنیم $f$ خطی است:
\[
f(x) = w^T x
\]
(جمله بایاس را می‌توان با افزودن یک مقدار ثابت $1$ به $x$ در $w$ ادغام کرد.)

پارامترهای $w$ مجموع خطاهای مربعات (\lr{sum of squared errors}) را کمینه می‌کنند:
\[
\min_{w} \sum_{i=1}^{n} (y_i - w^T x_i)^2
\]

\section*{۲. رگرسیون هسته: ایده اصلی}

\subsection*{چرا هسته‌ها؟}
داده‌های دنیای واقعی اغلب دارای الگوهای غیرخطی هستند که مدل‌های خطی نمی‌توانند آنها را ثبت کنند. راه حل: نگاشت هر ورودی $x$ به یک فضای ویژگی با ابعاد بالاتر $\phi(x)$ است که در آن رابطه خطی می‌شود:
\[
f(x) = w^T \phi(x)
\]

\textbf{چالش:} محاسبه صریح $\phi(x)$ ممکن است از نظر محاسباتی بسیار سنگین یا حتی دارای ابعاد نامتناهی باشد.

\textbf{ترفند هسته (\lr{Kernel Trick}):} استفاده از یک تابع هسته که ضرب‌های داخلی را در فضای ویژگی بدون ساخت صریح $\phi(x)$ محاسبه می‌کند:
\[
K(x_i, x_j) = \langle \phi(x_i), \phi(x_j) \rangle
\]

\subsection*{فرمولاسیون دوگان (\lr{Dual Formulation})}
با استفاده از ترفند هسته، تابع رگرسیون به مجموع وزنی شباهت‌های هسته با هر نقطه آموزشی تبدیل می‌شود:
\[
\hat{y}(x) = \sum_{i=1}^{n} \alpha_i K(x, x_i)
\]
که در آن:
\begin{itemize}
    \item $\alpha_i$ وزن‌های یادگرفته‌شده از داده‌ها هستند.
    \item $K(x, x_i)$ شباهت بین نقطه تست $x$ و نقطه آموزشی $x_i$ را اندازه‌گیری می‌کند.
\end{itemize}

\textbf{تفسیر:} پیش‌بینی یک میانگین وزن‌دهی‌شده بر اساس شباهت خروجی‌های آموزشی است، که در آن وزن‌های $\alpha_i$ یاد گرفته می‌شوند و میانگین‌های ساده نیستند.

\section*{۳. آموزش: حل برای $\alpha$}

\subsection*{مسئله بهینه‌سازی}
ما $\alpha$ را با برازش داده‌های آموزشی یاد می‌گیریم:
\[
\min_{\alpha} \sum_{j=1}^{n} \left( y_j - \sum_{i=1}^{n} \alpha_i K(x_j, x_i) \right)^2
\]

\subsection*{فرمولاسیون ماتریسی}
تعریف کنید:
\begin{itemize}
    \item ماتریس هسته (\lr{Kernel matrix}) (اندازه $n \times n$): $K_{ij} = K(x_i, x_j)$
    \item بردار وزن: $\alpha = [\alpha_1, \dots, \alpha_n]^T$
    \item بردار خروجی: $y = [y_1, \dots, y_n]^T$
\end{itemize}
برای تمام نقاط آموزشی، پیش‌بینی‌ها عبارتند از: $\hat{y} = K \alpha$. مجموع خطاهای مربعات به صورت زیر در می‌آید: $\|y - K \alpha\|^2$.

\subsection*{راه حل}
گرفتن مشتق و قرار دادن آن برابر صفر:
\[
\frac{\partial}{\partial \alpha} \|y - K \alpha\|^2 = -2 K^T(y - K \alpha) = 0
\]
از آنجا که $K$ متقارن است ($K^T = K$):
\[
K(y - K \alpha) = 0 \implies K y = K^2 \alpha
\]
اگر $K$ وارون‌پذیر باشد (هسته معین مثبت با نقاط متمایز)، در $K^{-1}$ ضرب می‌کنیم:
\[
y = K \alpha \implies \alpha = K^{-1} y
\]

\textbf{بینش کلیدی:} وزن‌های $\alpha$ به گونه‌ای انتخاب می‌شوند که پیش‌بینی‌های آموزشی کاملاً با خروجی‌های آموزشی مطابقت داشته باشند ($\hat{y}_{\text{train}} = y$). این درون‌یابی دقیق (\lr{exact interpolation}) است.

\section*{۴. پیش‌بینی روی داده‌های جدید}

\subsection*{تک نقطه تست}
برای یک نقطه تست جدید $x_{\text{test}}$ (که در مجموعه آموزشی نیست):
\begin{itemize}
    \item \textbf{مرحله ۱:} محاسبه شباهت‌های هسته بین $x_{\text{test}}$ و تمام نقاط آموزشی:
    \[
    k_{\text{test}} = [K(x_{\text{test}}, x_1), K(x_{\text{test}}, x_2), \dots, K(x_{\text{test}}, x_n)]^T \quad (\text{اندازه } n \times 1)
    \]
    \item \textbf{مرحله ۲:} انجام پیش‌بینی با استفاده از همان فرمول:
    \[
    \hat{y}_{\text{test}} = \sum_{i=1}^{n} \alpha_i K(x_{\text{test}}, x_i) = k_{\text{test}}^T \alpha
    \]
\end{itemize}

\subsection*{نقاط تست متعدد}
برای $m$ نقطه تست $\{x^{(1)}_{\text{test}}, \dots, x^{(m)}_{\text{test}}\}$:
ماتریس $K_{\text{test}}$ را با اندازه $m \times n$ تعریف کنید:
\[
[K_{\text{test}}]_{pj} = K(x^{(p)}_{\text{test}}, x_j)
\]
سپس پیش‌بینی‌ها:
\[
\hat{y}_{\text{test}} = K_{\text{test}} \alpha
\]

\subsection*{خلاصه آموزش در مقابل تست}
مهم: ویژگی درون‌یابی $\hat{y}_{\text{train}} = y$ فقط برای داده‌های آموزشی برقرار است. برای نقاط تست، پیش‌بینی‌ها بر اساس شباهت به نقاط آموزشی تعیین می‌شوند، نه مطابقت دقیق.
\begin{itemize}
    \item \textbf{آموزش:} ورودی $K_{n \times n}$، پیش‌بینی $\hat{y}_{\text{train}} = K \alpha$، برابر با مقدار واقعی؟ بله (برازش دقیق)
    \item \textbf{تست:} ورودی $k_{\text{test}}$ با اندازه $n \times 1$، پیش‌بینی $\hat{y}_{\text{test}} = k_{\text{test}}^T \alpha$، برابر با مقدار واقعی؟ خیر (تعمیم)
\end{itemize}

\section*{۵. ارتباط با شبه‌معکوس}
به یاد آورید که شبه‌معکوس $X^+$ معادله $X w = y$ را به طور بهینه حل می‌کند. برای رگرسیون هسته، ما $K \alpha = y$ را حل می‌کنیم، بنابراین:
\[
\alpha = K^+ y
\]
وقتی $K$ وارون‌پذیر باشد، $K^+ = K^{-1}$. وقتی $K$ منفرد یا بدحالت باشد، از منظم‌سازی (\lr{regularization}) استفاده می‌کنیم.

\subsection*{منظم‌سازی (رگرسیون ستیغی هسته - \lr{Kernel Ridge Regression})}
در عمل، با داده‌های نویزدار، درون‌یابی دقیق منجر به بیش‌برازش (\lr{overfitting}) می‌شود. یک جریمه ستیغی $\lambda \|f\|^2$ اضافه کنید:
\[
\min_{\alpha} \|y - K \alpha\|^2 + \lambda \alpha^T K \alpha
\]
راه حل به صورت زیر در می‌آید:
\[
\alpha = (K + \lambda I)^{-1} y
\]

این تخمین‌گر رگرسیون ستیغی هسته است. اصطلاح منظم‌سازی:
\begin{itemize}
    \item پایداری عددی (\lr{numerical stability}) را بهبود می‌بخشد (وارون‌پذیری $K + \lambda I$ را تضمین می‌کند).
    \item با کوچک کردن وزن‌ها از بیش‌برازش جلوگیری می‌کند.
    \item همواری تابع پیش‌بینی را کنترل می‌کند.
\end{itemize}

\section*{۶. تاثیر منظم‌سازی و سطح نویز در رگرسیون هسته}
\begin{itemize}
    \item \textbf{نویز $\sigma = 0$ (بدون نویز):}
    \begin{itemize}
        \item $\lambda = 0$: برازش کامل ($\text{MSE} \approx 0$) - مدل دقیقاً درون‌یابی می‌کند.
        \item $\lambda > 0$: انحراف جزئی ایجاد می‌شود، اما همچنان برازش خوبی است.
    \end{itemize}
    \item \textbf{نویز $\sigma = 0.1$ (نویز کم):}
    \begin{itemize}
        \item $\lambda = 0$: ممکن است تا حدی بیش‌برازش رخ دهد و الگوهای نویز را ثبت کند.
        \item $\lambda = 10^{-6}$: تعادلی بین برازش و همواری برقرار می‌کند.
        \item $\lambda = 0.1$: ممکن است کم‌برازش (\lr{underfit}) رخ دهد (بسیار هموار).
    \end{itemize}
    \item \textbf{بینش‌های کلیدی:}
    \begin{itemize}
        \item نویز بیشتر $\to$ نیاز به منظم‌سازی قوی‌تر ($\lambda$ بزرگتر).
        \item بدون نویز $\to$ $\lambda=0$ درون‌یابی دقیق می‌دهد.
        \item منظم‌سازی بیش از حد $\to$ کم‌برازش (انحراف یا \lr{bias}).
        \item منظم‌سازی بسیار کم $\to$ بیش‌برازش (واریانس).
    \end{itemize}
\end{itemize}

\section*{۷. دیدگاه نخستین (\lr{Primal}) در مقابل دوگان (\lr{Dual})}

\subsection*{دو شکل راه حل}
در رگرسیون خطی استاندارد با ویژگی‌های $X$، بردار وزن را می‌توان به دو روش معادل بیان کرد:
\begin{itemize}
    \item \textbf{فرم نخستین (\lr{Primal Form}) (فضای ویژگی):}
    \[
    w = (X^T X)^{-1} X^T y
    \]
    \item \textbf{فرم دوگان (\lr{Dual Form}) (فضای نمونه):}
    \[
    w = X^T (X X^T)^{-1} y
    \]
\end{itemize}
هر دو در صورتی که $X$ دارای رتبه ستونی کامل باشد، پیش‌بینی‌های یکسانی را ارائه می‌دهند.

\subsection*{چرا فرم دوگان برای هسته‌ها اهمیت دارد}
فرم دوگان نشان می‌دهد که $w$ یک ترکیب خطی از نمونه‌های آموزشی است:
\[
w = \sum_{i=1}^{n} \alpha_i x_i \quad \text{که در آن} \quad \alpha = (X X^T)^{-1} y
\]
وقتی به فضای ویژگی $\phi(x)$ نگاشت می‌کنیم، این به صورت زیر در می‌آید:
\[
w = \sum_{i=1}^{n} \alpha_i \phi(x_i) \quad \text{با} \quad \alpha = K^{-1} y
\]
که در آن $K_{ij} = \phi(x_i)^T \phi(x_j) = K(x_i, x_j)$ است.

\textbf{بینش کلیدی:} پیش‌بینی برای یک نقطه جدید به صورت زیر در می‌آید:
\[
\hat{y}(x) = w^T \phi(x) = \sum_{i=1}^{n} \alpha_i \phi(x_i)^T \phi(x) = \sum_{i=1}^{n} \alpha_i K(x_i, x)
\]
ما هرگز به طور صریح به $\phi(x)$ نیاز نداریم — فقط ارزیابی‌های هسته $K(x_i, x)$ مورد نیاز است!

\subsection*{موازنه‌های محاسباتی}
\begin{itemize}
    \item \textbf{سناریوی $n \ll d$ (نمونه‌های کم، ویژگی‌های زیاد):} فرم نخستین کند (معکوس کردن $d \times d$) و فرم دوگان سریع (معکوس کردن $n \times n$) است.
    \item \textbf{سناریوی $n \gg d$ (نمونه‌های زیاد، ویژگی‌های کم):} فرم نخستین سریع (معکوس کردن $d \times d$) و فرم دوگان کند (معکوس کردن $n \times n$) است.
    \item \textbf{هسته‌های غیرخطی (RBF، چندجمله‌ای):} فرم نخستین غیرقابل اعمال و فرم دوگان مورد نیاز است (ترفند هسته).
\end{itemize}

مثال: هسته \lr{RBF} با $n=1000$ نمونه به طور موثر از ویژگی‌های بی‌نهایت بعدی ($d=\infty$) استفاده می‌کند:
\begin{itemize}
    \item فرم نخستین غیرممکن است (نمی‌توان $\phi(x)$ را صراحتاً محاسبه کرد).
    \item فرم دوگان با ماتریس هسته $1000 \times 1000$ کار می‌کند.
\end{itemize}

\section*{۸. نکات کلیدی}
\begin{enumerate}
    \item رگرسیون هسته برازش غیرخطی را با نگاشت ضمنی داده‌ها به فضاهای با ابعاد بالا امکان‌پذیر می‌کند.
    \item فرمولاسیون دوگان پیش‌بینی‌ها را به عنوان مجموع‌های وزنی از شباهت‌های هسته بیان می‌کند: $\hat{y}(x) = \sum_{i=1}^{n} \alpha_i K(x, x_i)$.
    \item آموزش معادله $\alpha = K^{-1} y$ (بدون منظم‌سازی) یا $\alpha = (K + \lambda I)^{-1} y$ (منظم‌شده) را حل می‌کند.
    \item تست از فرمول $\hat{y}_{\text{test}} = k_{\text{test}}^T \alpha$ استفاده می‌کند، که در آن $k_{\text{test}}$ شامل شباهت‌های بین نقاط تست و آموزشی است.
    \item منظم‌سازی ($\lambda > 0$) برای موارد زیر ضروری است:
    \begin{itemize}
        \item پایداری عددی (تضمین اینکه $K + \lambda I$ وارون‌پذیر است).
        \item جلوگیری از بیش‌برازش به نویز.
    \end{itemize}
    \item انتخاب هسته انعطاف‌پذیری مدل را کنترل می‌کند:
    \begin{itemize}
        \item هسته خطی $\to$ رگرسیون خطی استاندارد.
        \item هسته \lr{RBF} $\to$ تقریب هموار و محلی.
        \item هسته چندجمله‌ای $\to$ جملات متقابل تا درجه $d$.
    \end{itemize}
\end{enumerate}

\end{document}
